\section{Background}
\label{sec:background}

Performance of web services is sensitive to system metrics of their backend servers, in container-based cluster, that is backend containers. For convenience, we called the backend containers as endpoints in this paper. This section first points out factors that affect performance of web services in container-based cluster, and then introduces the current design of load-balancing solution, that is ingress controller in Kubernetes and finally state defects of ingress controller.

\subsection{Influencing factors}
\label{subsec:Influencing_factors}

Factors that affect the performance include CPU usage, memory usage, network usage, services dependency and locality.
As known to us, servers with high usage in CPU, memory or network will perform badly. Things are same when it comes to container-based cluster, for an endpoint, it causes delay on requests processing when it is high-loaded in CPU and network. More and more, if the endpoints are in high-loaded state for memory, they face the risk that they may be restarted due to OOM(out of memory) problem. Thus, if we demand high performance for web application, we'd better even load of each endpoint and reduce overhead endpoints as many as possible.
services dependency means that services do not provides functions themselves, they rely on applications of other services to finish their jobs. For example, we have endpoint1, endpoint2,endpoint3 for serviceA, and endpoint4,enpoint5 for serviceB, and serviceA relies on serviceB to process its requests, particularly, endpoint1 and endpoint2 rely on endpoint5, endpoint3 relies on endpoint5. Obviously, the performance of endpoint4 and endpoint5 will directly influence the performance of endpoint1, endpoint2 and endpoint3. Furthermore, the competition between endpoint1 and endpoint2 for accessing endpoint5 will have a negatice impact on performance of serviceA. This paper select database dependency as an example of services dependency. That is endpoints of a web service rely on endpoints of database to process requests, which is also a common case for web services. Traditionally, developers tend to reduce pressure of database by dividing it into several databases and synchronize data between them. Therefore, it is common for an endpoint to access a database that is not only the same as some endpoints, but also different from others. Under this circumstance, performance of web applications can be improved through reducing race between endpoints that share the same database or reducing the overhead endpoints of databases.
Locality means that the placement of endpoints may affect their performance. Though containers are claimed to isolate resource from each other, papers~{\cite{Zhao2018Locality}} still show that when two or more endpoints are placed in the same node, the delay appears. More and more, although under circumstance that an endpoint monopolizes a node, it still remains possibility for delay when the node is in high-load status because of some system processes. So the locality should also be considered for a better load-balancing performance.

\subsection{Ingress controller}
\label{subsec:ingress_controller}

\begin{figure}[htbp]
  \centering
  \includegraphics[width=0.45\textwidth]{images/ingress_controller.png}\\
  \caption{Workflow of ingress controller}
  \label{fig:ingress_controller}
\end{figure}
\hspace{0pt}

Figure~{\ref{fig:ingress_controller}} shows the workflow of kubernetes load-balancing component -- ingress controller: To balance network workload, ingress controller focuses on two part, one is to keep track of the IP and number change of endpoints so that it can update traffic rules in time to guarantee accuracy. The other part is watching ingress to update load-balancing strategy. Ingress is a resource describing inbound rules for connections to reach endpoints. With correct status of endpoints and a configured ingress, ingress controller sets up appropriate traffic rules for its built-in loadbalancer(nginx). Finally, the loadbalancer dispatches requests to suitable endpoints according to configured rules, balancing load across endpoints.



\subsection{Shortcomings of ingress controller}
\label{subsec:Shortcoming}
Rather than balancing load, ingress controller concentrates more on linking ingress to nginx features like session affinity to let users to custom their load-balancing algorithm smoothly in kubernetes. However, its configurable algorithms are simply based on static algorithms like round\_robin, which do not consider influencing factors mentioned above. To summary, the igress controller has two shortcomings when it is put into pratice, that is hysteresis and one-sidedness. The round\_robin strategy assumes that all requests have the same resource consumption and aims to balance load after a long period of time, so problems come when requests change in a rather short time. For example, when there is a peak of requests, and new endpoints is needed to participate in providing service. A appropriate way to balance load now is to dispatch more workload to new idle endpoints and reduce workload for high-loaded old endpoints, but the round\_robin algorithm still dispatches requests indistinguishably, taking the riks that the old high-loaded endpoints may be unable to work normally, even restarted due to OOM although it will finally balance load relatively. The one-sideness means that the round\_robin algorithm only considers to even workload accross endpoints, without taking influencing factors mentioned above, thus failes to balance load accurately. To address these shortcomings of ingress controller, this paper presents an alternative basic algorithm to ingress controller, that is a dynamic load algorithm that takes these factors and request information into consideration.
